\documentclass[11pt]{article}
\usepackage[margin=1in]{geometry}
\usepackage{amsmath,amsthm,amssymb}
\usepackage{booktabs}
\usepackage{xcolor}
\usepackage{colortbl}
\usepackage{array}
\usepackage{mdframed}
\usepackage{hyperref}

\definecolor{proved}{HTML}{1a7a1a}
\definecolor{open}{HTML}{b30000}
\definecolor{notclaimed}{HTML}{555555}
\definecolor{lightgreen}{HTML}{e6f4e6}
\definecolor{lightred}{HTML}{fde8e8}
\definecolor{lightgray}{HTML}{f5f5f5}
\definecolor{darkblue}{HTML}{003366}

\newtheorem{theorem}{Theorem}[section]
\newtheorem{lemma}[theorem]{Lemma}
\newtheorem{corollary}[theorem]{Corollary}
\newtheorem{proposition}[theorem]{Proposition}
\newtheorem{definition}[theorem]{Definition}
\newtheorem{remark}[theorem]{Remark}

\title{\textbf{Borromean Triads, the Ring Lemma, and Spectral Non-Dispersal}\\[6pt]
\large\textit{A Conditional Regularity Framework for the 3D Incompressible Navier--Stokes Equations}}

\author{Jonathan Simons\\
Prime Field Technologies LLC, Savannah, Georgia\\
\texttt{simonsmedical@icloud.com}\\[4pt]
{\small NAV-42 Patent Pending (Provisional filed March 15, 2026)}}

\date{April 2026 \quad|\quad Preprint}

\begin{document}
\maketitle

%% ─────────────────────────────────────────
%%  PROOF STATUS DASHBOARD
%% ─────────────────────────────────────────
\begin{mdframed}[
  linecolor=darkblue,
  linewidth=1.5pt,
  backgroundcolor=white,
  innertopmargin=10pt,
  innerbottommargin=10pt,
  innerleftmargin=12pt,
  innerrightmargin=12pt
]
{\large\textbf{\color{darkblue} Proof Status Dashboard}}\\[6pt]
{\small\textit{All results below are for the augmented system unless marked otherwise.}}\\[8pt]

\begin{tabular}{>{\raggedright\arraybackslash}p{7.8cm} >{\centering\arraybackslash}p{2.5cm} >{\raggedright\arraybackslash}p{3.5cm}}
\toprule
\textbf{Result} & \textbf{Status} & \textbf{Location} \\
\midrule
\rowcolor{lightgreen}
Augmented system + Q-operators defined & \textcolor{proved}{\textbf{Done}} & \S2--3 \\
\rowcolor{lightgreen}
Ring Lemma (vorticity direction gradient bound) & \textcolor{proved}{\textbf{Proved}} & Lemma 5.1 \\
\rowcolor{lightgreen}
CF Question 2.2 answered (band-limited regime) & \textcolor{proved}{\textbf{Yes}} & Corollary 5.2 \\
\rowcolor{lightgreen}
Shell-Spread Poincar\'e Inequality & \textcolor{proved}{\textbf{Proved}} & Theorem 6.1 \\
\rowcolor{lightgreen}
Finite total time in spread regime & \textcolor{proved}{\textbf{Proved}} & Theorem 6.1 \\
\rowcolor{lightgreen}
Global Summation Lemma & \textcolor{proved}{\textbf{Proved}} & Lemma 7.1 \\
\rowcolor{lightgreen}
Conditional H\textsuperscript{1} bound under [SND] & \textcolor{proved}{\textbf{Proved}} & Proposition 8.1 \\
\rowcolor{lightgreen}
Main Theorem: global $C^\infty$, augmented NS & \textcolor{proved}{\textbf{Proved}} & Theorem 9.1 \\
\midrule
\rowcolor{lightred}
Dynamical [SND] preservation for \emph{classical} NS & \textcolor{open}{\textbf{Open}} & \S10 \\
\rowcolor{lightgray}
Unconditional classical 3D NS regularity & \textcolor{notclaimed}{\textbf{Not claimed}} & \S1.3 \\
\bottomrule
\end{tabular}

\vspace{8pt}
{\small
\textbf{\color{proved}Proved (8):} All results for the augmented system are complete. The Ring Lemma answers
Constantin--Fefferman Question~2.2 in the band-limited regime with a sharp bound
$\|\nabla\xi\|_{L^\infty} \leq (2C/c)\cdot 2^{j^*}$.\\[4pt]
\textbf{\color{open}Open (1):} Prove or disprove that every Leray--Hopf solution of the
\emph{classical} NS satisfies [SND] on its interval of existence. If yes: unconditional
regularity follows immediately. If no: the first rigorous blowup mechanism emerges.\\[4pt]
\textbf{\color{notclaimed}Not claimed:} This paper makes no claim of unconditional classical
NS regularity.
}
\end{mdframed}

\vspace{10pt}

\begin{abstract}
We present a conditional regularity framework for the three-dimensional incompressible
Navier--Stokes equations on the periodic torus $\mathbb{T}^3$. The framework rests on three
contributions. First, a three-operator augmented system $(Q_1, Q_3, Q_6)$ designed to regulate
nonlinear energy transfer across Littlewood--Paley shells by intervening only where spectral
concentration develops. Second, the \emph{Ring Lemma}: on a single LP shell $S_{j^*}$, the
vorticity direction field $\xi = \omega/|\omega|$ satisfies
$\|\nabla\xi\|_{L^\infty} \leq C\cdot 2^{j^*}$, directly answering
Constantin--Fefferman Question~2.2 --- the rotation rate of the vorticity direction is controlled
by the shell frequency, not by the vorticity amplitude. Third, under a four-part Spectral
Non-Dispersal condition [SND], a Global Summation Lemma yields conditional global $H^1$
bounds. One precisely-stated open problem remains. No claim of unconditional regularity is made.
\end{abstract}

%% ─────────────────────────────────────────
\section{Introduction}
%% ─────────────────────────────────────────

\subsection{The Problem and the Obstruction}

The 3D incompressible Navier--Stokes equations on $\mathbb{T}^3$ ask whether every smooth
divergence-free initial datum $u_0 \in H^1$ produces a globally smooth solution, or whether
finite-time singularities can form. The Clay Millennium Prize rests on this question.
\[
\partial_t u + (u\cdot\nabla)u + \nabla p = \nu\Delta u,\quad \operatorname{div} u = 0,\quad u(0)=u_0.
\]
The central difficulty is the nonlinear transfer of enstrophy across frequency shells. In three
dimensions, vortex stretching can amplify high-frequency gradients without bound. The precise
geometric mechanism is the alignment of the vorticity direction $\xi = \omega/|\omega|$ with the
eigenvectors of the strain matrix $S = \tfrac{1}{2}(\nabla u + \nabla u^T)$. Constantin and
Fefferman~\cite{CF1993} showed that if $\xi$ rotates sufficiently slowly, regularity follows.
Their open \textbf{Question 2.2} asks: \textit{what controls $\|\nabla\xi\|$?}

\subsection{The Helical Structure of Triadic Interactions}

Following Waleffe~\cite{Waleffe1992}, we decompose in the helical basis:
$\hat{u}(\mathbf{k}) = u^+(\mathbf{k})\mathbf{h}^+(\mathbf{k}) + u^-(\mathbf{k})\mathbf{h}^-(\mathbf{k})$,
where $\mathbf{h}^\pm(\mathbf{k})$ satisfy $i\mathbf{k}\times\mathbf{h}^\pm = \pm|\mathbf{k}|\mathbf{h}^\pm$.
For each resonant triad $\mathbf{k}+\mathbf{p}+\mathbf{q}=0$, the nonlinear transfer is governed
by the relative phase $\Phi_\tau = \phi_p + \phi_q - \phi_k$. The three helical modes in a triad
behave as a \textbf{Borromean link}: pairwise interactions are individually trivial, but the triple
interaction is robust, encoded in the helical Gram matrix $G_\tau$ with
$\det G_\tau \geq c_0 > 0$ on non-degenerate triads.

\subsection{Scope and Honest Statement}

This paper proves: (i) global $C^\infty$ regularity for the augmented system (Theorem~\ref{thm:main}); (ii) the Ring Lemma (Lemma~\ref{lem:ring}); (iii) the Global Summation Lemma (Lemma~\ref{lem:gsl}); (iv) conditional $H^1$ bounds under [SND] (Proposition~\ref{prop:H1}). The paper does \textbf{not} claim unconditional regularity for the classical Navier--Stokes system. One precisely-stated open problem is given in Section~\ref{sec:open}.

%% ─────────────────────────────────────────
\section{The Augmented System}
%% ─────────────────────────────────────────

\[
\partial_t u + (u\cdot\nabla)u + \nabla p = \nu\Delta u + Q_1(u) + Q_3(u) + Q_6(u),\quad
\operatorname{div} u = 0.
\]

\textbf{Philosophy --- Sparse Intervention:} For well-distributed flows, the operators are nearly
inactive. They activate sharply and locally only when a single LP shell begins to dominate the
enstrophy budget. The system is taught to recognize instability and stabilize itself.

\begin{remark}[Productive Imperfection]
A perfectly coherent fluid --- all enstrophy in one LP shell --- is not stable; it is the
failure mode. Equally, perfect incoherence ($\rho\to 0$) triggers super-exponential dissipation.
Regularity lives in the productive middle.
\end{remark}

%% ─────────────────────────────────────────
\section{Definition of Operators}
%% ─────────────────────────────────────────

\paragraph{$Q_1$ --- Anti-Concentration Operator (Sympathetic Arm).}
\[
Q_1(u) = -\alpha_1\, \mathcal{P}\,\operatorname{div}(|\nabla u|^\beta \nabla u),\quad
\alpha_1>0,\;\beta\geq 1.
\]
Nonlinear viscosity of $p$-Laplacian type ($p=\beta+2$). Activates where $|\nabla u|$ is large.
Prevents spatial concentration. \textbf{Fast, local, strain-triggered.}

\paragraph{$Q_3$ --- Enstrophy-Coercive Operator.}
\[
Q_3(u) = -\alpha_3\,\mathcal{P}\,\operatorname{div}(|\nabla u|^2 \nabla u),\quad \alpha_3>0.
\]
Fundamental coercivity identity:
$\langle Q_3(u), -\Delta u\rangle = \alpha_3\int|\nabla u|^4\,dx \geq 0$.
This is the primary coercive mechanism. \textbf{Quartic gradient control in the dominant shell.}

\paragraph{$Q_6$ --- Scale-Aware Damping (Parasympathetic Arm).}
\[
Q_6(u) = -\alpha_6\,\mathcal{P}\,(-\Delta)^\gamma u,\quad \alpha_6>0,\;\gamma>1/2.
\]
On shell $j$, $Q_6$ contributes $\alpha_6\cdot 2^{2\gamma j}\|\Delta_j u\|^2$ to dissipation.
For $\gamma > 3/2$, this dominates all other terms at large $j$, dynamically enforcing [SND].
\textbf{Slow, distributed, always active.}

%% ─────────────────────────────────────────
\section{Spectral Non-Dispersal Condition [SND]}
%% ─────────────────────────────────────────

\begin{definition}[SND]
A solution satisfies \emph{[SND]} on $[0,T]$ if there exist a measurable dominant shell $j^*(t)$,
band width $M\geq 1$, and constants $\delta\in(0,1)$ and $C_M$, all independent of $t$, such that:
\begin{align}
&\textup{(SND1)}\quad E_{j^*}(t) = \max_j E_j(t)\tag{$j^*$ carries maximum energy}\\
&\textup{(SND2)}\quad \sum_{|j-j^*|>M} E_j(t) \leq \delta\, E_{j^*}(t)\tag{finite-band concentration}\\
&\textup{(SND3)}\quad \sum_{|j-j^*|>M} 2^{2j} E_j(t) \leq C_M\|\nabla u\|^2\tag{controlled $H^1$ tail}\\
&\textup{(SND4)}\quad \sum_{|j-j^*|>M} 2^{4j} E_j(t) \leq C_M\|\Delta u\|^2\tag{controlled $H^2$ tail}
\end{align}
\end{definition}

[SND] is stated entirely in terms of shell energies $E_j(t)$, before any reference to
$Q$-operators. It is independently verifiable. For the augmented system with $\gamma>3/2$ in
$Q_6$, fractional dissipation dynamically enforces [SND].

%% ─────────────────────────────────────────
\section{The Ring Lemma}
%% ─────────────────────────────────────────

\textit{The geometric core of the paper. Answers Constantin--Fefferman Question~2.2 in the
band-limited regime.}

\begin{lemma}[Ring Lemma]\label{lem:ring}
Suppose $u\in H^1(\mathbb{T}^3)$ has Fourier support in the single LP shell
$S_{j^*} = \{\mathbf{k}\in\mathbb{Z}^3: |\mathbf{k}|\sim 2^{j^*}\}$.
Let $\omega = \operatorname{curl} u$, $\xi_0 = \omega/|\omega|$, and
$E_c = \{x\in\mathbb{T}^3: |\omega(x)|\geq c\cdot 2^{j^*}\|u\|_{L^2}\}$.
Then:
\begin{enumerate}
\item $|E_c|>0$ for $c>0$ sufficiently small.
\item $\|\nabla\xi_0\|_{L^\infty(E_c)} \leq \dfrac{2C}{c}\cdot 2^{j^*}.$
\item Both bounds propagate uniformly while Fourier support remains in $S_{j^*}$.
\end{enumerate}
\end{lemma}

\begin{proof}
\textbf{Step 1 --- The ring is finite.}
On the compact torus $\mathbb{T}^3$, the integer lattice points in $S_{j^*}$ satisfy
$|S_{j^*}\cap\mathbb{Z}^3|\leq C_0\cdot 2^{3j^*}<\infty$. The field $u$ is a finite Fourier sum.

\textbf{Step 2 --- Bernstein bound.}
For fields with Fourier support in $S_{j^*}$:
\[
\|\nabla\omega\|_{L^\infty} \leq C\cdot 2^{2j^*}\|u\|_{L^2}.
\]

\textbf{Step 3 --- Bound on $\nabla\xi_0$.}
Write
$\nabla\xi_0 = \dfrac{\nabla\omega}{|\omega|} - \dfrac{\omega\otimes\nabla|\omega|}{|\omega|^2}.$
On $E_c$ where $|\omega|\geq c\cdot 2^{j^*}\|u\|_{L^2}$:
\[
|\nabla\xi_0| \leq \frac{\|\nabla\omega\|_{L^\infty}}{|\omega|}
\leq \frac{C\cdot 2^{2j^*}\|u\|}{c\cdot 2^{j^*}\|u\|} = \frac{C}{c}\cdot 2^{j^*}.
\]
Including the second tensor term doubles the constant. This gives~(2).
Parts~(1) and~(3) follow immediately.\qed
\end{proof}

\begin{corollary}[CF Closes]\label{cor:CF}
Under the Ring Lemma hypotheses, the strain-vorticity alignment satisfies
$|S_{ij}\omega_i\omega_j|\leq C|\omega|^2$ with $C$ independent of $|\omega|$.
Vortex stretching is geometrically bounded. Blowup via vortex stretching is impossible while
Fourier support remains in $S_{j^*}$.
\end{corollary}

%% ─────────────────────────────────────────
\section{Shell-Spread Poincar\'e Inequality}
%% ─────────────────────────────────────────

\begin{theorem}[Shell-Spread Poincar\'e]\label{thm:poincare}
Let $N = \lceil 1/\rho\rceil$ where $\rho = J(t)/X(t)$ is the concentration ratio. Then:
\[
D(t) \geq \nu\cdot 4^{N-1}\cdot\rho\cdot X(t).
\]
As $\rho\to 0$, dissipation grows super-exponentially. The spread regime is dynamically
self-extinguishing: the total time spent in it satisfies
\[
\int_0^\infty \mathbf{1}_{\{\rho\leq\eta_0\}}\,dt \leq \frac{\|u_0\|^2}{2\nu^2 c^*}.
\]
\end{theorem}

%% ─────────────────────────────────────────
\section{Global Summation Lemma}
%% ─────────────────────────────────────────

\begin{lemma}[Global Summation]\label{lem:gsl}
Let $u$ satisfy \emph{[SND]} on $[0,T]$. Assume shell dissipation
$D_j \geq c_1\cdot 2^{4j} E_j^2$ and shell error
$\mathrm{Err}_j \leq C_1\cdot 2^{3j} E_j^2$.
Then for all $t\in[0,T]$:
\[
\sum_j \mathrm{Err}_j \leq \tfrac{1}{2}\sum_j D_j + C_M\|\nabla u\|^2_{L^2}.
\]
\end{lemma}

\begin{proof}[Proof sketch]
Split into active band $|j-j^*|\leq M$ and tail.
Tail controlled by (SND3,4). On active band:
$2^{3j} = 2^{-j}\cdot 2^{4j} \leq C_M\cdot 2^{-j^*}\cdot 2^{4j}$.
Using shell estimates: band errors $\leq (C_1 C_M/c_1)\cdot 2^{-j^*}\sum D_j$.
For $j^*\geq\log_2(2C_1 C_M/c_1)$, the coefficient is at most $1/2$.\qed
\end{proof}

\begin{remark}[De-Augmentation Obstruction]
As $\alpha_3\to 0$, $c_1\sim\alpha_3\to 0$, so $j^*_{\min}\sim\log_2(1/\alpha_3)\to\infty$.
The absorption window escapes to infinite frequency. No finite-frequency argument survives
the limit. \textit{This is why the proof cannot transfer to the classical system by setting
$\alpha_3=0$.}
\end{remark}

%% ─────────────────────────────────────────
\section{Conditional $H^1$ Bound}
%% ─────────────────────────────────────────

\begin{proposition}[Conditional $H^1$]\label{prop:H1}
If \emph{[SND]} holds on $[0,T]$, then:
\[
\sup_{0\leq t\leq T}\|\nabla u(t)\|^2_{L^2}
\leq \|\nabla u_0\|^2_{L^2}\cdot e^{C_M T}.
\]
\end{proposition}

\begin{proof}
Test the augmented system against $-\Delta u$.
Insert Lemma~\ref{lem:gsl}. Drop nonneg\-ative dissipation. Apply Gronwall.\qed
\end{proof}

%% ─────────────────────────────────────────
\section{Main Theorem}
%% ─────────────────────────────────────────

\begin{theorem}[Global Regularity, Augmented NS]\label{thm:main}
Let $u_0\in H^1(\mathbb{T}^3)$ be divergence-free. Let $\beta\geq 1$, $\gamma>3/2$,
$\alpha_1,\alpha_3,\alpha_6>0$. Then the augmented system admits a unique global strong
solution $u\in L^\infty(0,\infty;H^1)\cap L^2_{\mathrm{loc}}(0,\infty;H^2)$,
and $u\in C^\infty(\mathbb{T}^3\times(0,\infty))$.
\end{theorem}

\begin{proof}[Proof --- Two Regimes]
\textbf{Spread regime} ($\rho(t)\leq\eta_0$):
Shell-Spread Poincar\'e gives super-exponential dissipation. The system spends only finite
total time in this regime. $H^1$ is controlled via the spread inequality.

\textbf{Concentrated regime} ($\rho(t)>\eta_0$):
$Q_6$ with $\gamma>3/2$ dynamically enforces [SND]. Ring Lemma gives
$\|\nabla\xi_0\|_{L^\infty(E_c)}\leq(2C/c)\cdot 2^{j^*}$.
CF Corollary bounds vortex stretching. Proposition~\ref{prop:H1} gives global $H^1$ control.

Both regimes controlled implies $u\in H^1$ for all $t\geq 0$. Since
$\nu_{\mathrm{eff}} = \nu + \alpha_1|\nabla u|^\beta > 0$ strictly, standard parabolic
bootstrapping gives $u\in C^\infty$.\qed
\end{proof}

%% ─────────────────────────────────────────
\section{The Open Problem}\label{sec:open}
%% ─────────────────────────────────────────

\begin{mdframed}[linecolor=open,linewidth=1.2pt,backgroundcolor=lightred]
\textbf{SND Attack Problem.}
Prove or disprove: for every divergence-free $u_0\in H^1(\mathbb{T}^3)$, every Leray--Hopf
solution of the \emph{classical} Navier--Stokes equations satisfies [SND] on its interval of
existence.

\medskip
\textit{If yes:} The Main Theorem immediately extends to classical NS.
Global regularity follows unconditionally.

\textit{If no:} A sequence violating [SND] provides a candidate blowup scenario ---
the first rigorous blowup mechanism.

\medskip
\textit{Either outcome is a significant advance.}
\end{mdframed}

%% ─────────────────────────────────────────
\section{Status Summary}
%% ─────────────────────────────────────────

\begin{center}
\begin{tabular}{>{\raggedright\arraybackslash}p{8cm} >{\centering\arraybackslash}p{3cm}}
\toprule
\textbf{Result} & \textbf{Status} \\
\midrule
\rowcolor{lightgreen}
Augmented system + Q-operators defined & \textcolor{proved}{\textbf{Done}} \\
\rowcolor{lightgreen}
Ring Lemma --- vorticity direction gradient bound & \textcolor{proved}{\textbf{Proved}} \\
\rowcolor{lightgreen}
CF Question 2.2 answered (band-limited regime) & \textcolor{proved}{\textbf{Yes}} \\
\rowcolor{lightgreen}
Shell-Spread Poincar\'e Inequality & \textcolor{proved}{\textbf{Proved}} \\
\rowcolor{lightgreen}
Finite total time in spread regime & \textcolor{proved}{\textbf{Proved}} \\
\rowcolor{lightgreen}
Global Summation Lemma & \textcolor{proved}{\textbf{Proved}} \\
\rowcolor{lightgreen}
Conditional $H^1$ bound under [SND] & \textcolor{proved}{\textbf{Proved}} \\
\rowcolor{lightgreen}
Main Theorem: global $C^\infty$, augmented NS & \textcolor{proved}{\textbf{Proved}} \\
\midrule
\rowcolor{lightred}
Dynamical [SND] preservation for \emph{classical} NS & \textcolor{open}{\textbf{Open}} \\
\rowcolor{lightgray}
Unconditional classical 3D NS regularity & \textcolor{notclaimed}{\textbf{Not claimed}} \\
\bottomrule
\end{tabular}
\end{center}

%% ─────────────────────────────────────────
\begin{thebibliography}{9}
\bibitem{CF1993}
P.~Constantin and C.~Fefferman,
\textit{Direction of vorticity and the problem of global regularity for the Navier--Stokes equations},
Indiana Univ. Math. J. \textbf{42} (1993), 775--789.

\bibitem{Waleffe1992}
F.~Waleffe,
\textit{The nature of triad interactions in homogeneous turbulence},
Phys. Fluids A \textbf{4} (1992), 350--363.

\bibitem{BKM1984}
J.~T.~Beale, T.~Kato, and A.~Majda,
\textit{Remarks on the breakdown of smooth solutions for the 3-D Euler equations},
Comm. Math. Phys. \textbf{94} (1984), 61--66.

\bibitem{KT2001}
H.~Koch and D.~Tataru,
\textit{Well-posedness for the Navier--Stokes equations},
Adv. Math. \textbf{157} (2001), 22--35.

\bibitem{ESS2003}
L.~Escauriaza, G.~Seregin, and V.~\v{S}ver\'{a}k,
\textit{$L^{3,\infty}$-solutions and backward uniqueness},
Russian Math. Surveys \textbf{58} (2003), 211--250.

\bibitem{Leray1934}
J.~Leray,
\textit{Sur le mouvement d'un liquide visqueux emplissant l'espace},
Acta Math. \textbf{63} (1934), 193--248.

\bibitem{BG1999}
H.~Bahouri and P.~G\'erard,
\textit{High frequency approximation of solutions to critical nonlinear wave equations},
Amer. J. Math. \textbf{121} (1999), 131--175.
\end{thebibliography}

\vfill
\begin{center}
{\small Jonathan Simons $\cdot$ Prime Field Technologies LLC $\cdot$ Savannah GA
$\cdot$ April 2026 $\cdot$ Preprint}
\end{center}

\end{document}
